\documentclass[10pt,a4paper]{amsart}
\usepackage[margin=19mm]{geometry}
\usepackage{amsmath,amssymb,mathtools}
\usepackage[scr=rsfs,cal=euler]{mathalfa}
\usepackage{xfrac}
\usepackage[unicode,hidelinks,bookmarks=false]{hyperref}
\newcommand{\ie}{i.e.\ }
\newcommand{\cf}{cf.\ }
\newcommand{\resp}{respectively\ }
\newcommand{\Subsection}{\subsection}
\providecommand{\nobreakdash}{\nobreak\hbox{-}}
\setlength{\emergencystretch}{2em}
\multlinegap0pt
\usepackage{bm}
\usepackage{bbm}
\usepackage{dsfont}
\usepackage{xspace}
\usepackage{cancel}
\usepackage{mathtools}
\usepackage{amsthm}
\makeatletter
\@ifundefined{theorem}{\newtheorem{theorem}{Theorem}}{}
\@ifundefined{lemma}{\newtheorem{lemma}[theorem]{Lemma}}{}
\makeatother
\makeatletter
\expandafter\ifx\csname customthm\endcsname\relax
\newenvironment{customthm}[1]
	{\renewcommand\theinnercustomthm{#1}\innercustomthm}
	{\endinnercustomthm}
\fi
\makeatother
\title[On the Rigidity of the Roots of Power Series with Constraine]{On the Rigidity of the Roots of Power Series with Constrained Coefficients}
\author{Jacob Kewarth}
\date{Source: arXiv:2505.18317v1 (CC BY 4.0)}
\begin{document}
\maketitle

\noindent\emph{Source and licence.} On the Rigidity of the Roots of Power Series with Constrained Coefficients. Version arXiv:2505.18317v1 (CC BY 4.0), distributed under the Creative Commons Attribution 4.0 International licence, as recorded in the arXiv licence metadata for this exact version and shown on the abstract page https://arxiv.org/abs/2505.18317v1. Full rights notice: licenses/arxiv-2505.18317-CC-BY-4.0.txt.\par\smallskip

\noindent\textit{Editorial scope.} Counted unit: the Lemma whose source label is \texttt{None} in the article (source0.tex lines 537--546, source label \texttt{None}), delivered together with the article's own complete original proof of it (source0.tex lines 548--605, one \texttt{proof} environment of 58 lines). No mathematics is written, repaired or re-derived: statement and proof are the article's own text line for line. Required context retained uncounted: none. Every definition, display and label of those lines is retained unchanged. Omitted from this package and not redistributed: 1--536, 547--547, 606--1164. Nothing inside a retained excerpt is omitted or altered: every retained body is delivered exactly as the article's lines are written, so no omission receipt of any kind accompanies this package. Citations: the retained excerpts carry no citation command at all, so no bibliography entry of the article is transcribed and no reference is redistributed. Reference adaptation: none was needed and none was made; every reference inside the delivered unit resolves inside it, so no unresolved reference or citation is produced. Printed slips kept literal: no printed word, symbol, hypothesis or number of the counted statement or of its proof was altered. Layout and class: the article's own class is \texttt{amsart} with packages including pdfpages, amssymb, amsmath, amsthm, thm-restate, graphicx, float; the wrapper is the lane's pinned \texttt{amsart} editorial wrapper at 10pt on A4, which restates the theorem-like environments the retained text needs and adds the article's own macro definitions that the retained text needs (none), with extra wrapper packages (bm, bbm, dsfont, xspace, cancel, mathtools). The delivered numbering is the unit's own. Assets: no retained span contains a figure environment, a graphics-inclusion command, a PDF-page inclusion, a table or a TikZ picture; no image file is copied into this package, the package declares no asset dependency and no image file is redistributed with it. Licence: version arXiv:2505.18317v1 of this article is distributed under the Creative Commons Attribution 4.0 International licence, as recorded in the arXiv licence metadata for this exact version and shown on the abstract page https://arxiv.org/abs/2505.18317v1. A full rights notice is delivered at licenses/arxiv-2505.18317-CC-BY-4.0.txt. Characters: every retained excerpt is pure ASCII, so the delivered engine, XeLaTeX with its default Latin Modern text font, reports no missing character.
\begin{lemma}
Let $S$ be a finite non-empty symmetric normalized set of real numbers and let $M = \max(S)$. Then:
\begin{enumerate}
\item If $P(z) = \sum_{j=0}^\infty p_jz^j$ is a power series with $p_j \in S$ and $p_0 \neq 0$ has two roots $\lambda_1, \lambda_2$ with $|\lambda_j| = \frac{1}{\sqrt{M}+1}$ and such that $\lambda_1 = \overline{\lambda_2}$ and $\lambda_1$ is in the first quadrant. Then $\lambda_1 = \lambda_2 = \frac{1}{\sqrt{M}+1}$ and (after possibly multiplying by -1) $P(z) = 1 - (2\sqrt{M}+1)z + \sum_{j=2}^\infty Mz^j$.

\item If $2\sqrt{M} + 1 \in S$ and if $P_n(z) = \sum_{j=0}^\infty p_{n, j}z^j$ are power series with roots $\lambda_{n, 1}$ and $\lambda_{n, 2}$ with $|\lambda_{n, j}| \rightarrow \frac{1}{\sqrt{M}+1}$ and $p_{n, j} \in S$, $p_{n, 0} = 1$, and $p_{n, 1} = -(2\sqrt{M}+1)$ then for $n$ large enough $\lambda_{n, j} \in \mathbb{R}$ and for every $m$ there is an $N$ such that $n > N$ implies that $p_{n, k} = p_k$ for all $k = 0, 1, ..., m$.

\item if $P_n(z) = \sum_{j=0}^\infty p_{n, j}z^j$ are power series with roots $\lambda_{n, 1}$ and $\lambda_{n, 2}$ with $|\lambda_{n, j}| \rightarrow \frac{1}{\sqrt{M}+1}$ and $p_{n, j} \in S$ then for $n$ large enough (and after possibly multiplying by -1) $p_{n, 0} = 1$ and $p_{n, 1} = -(2\sqrt{M}+1)$.
\end{enumerate}
\end{lemma}

\begin{proof}[Proof of (1)]

\

Note that for $P$ to have two roots so deep into the disc we must have $p_0 = \pm 1$. Possibly multiplying the series by $-1$ we may assume $p_0 = 1$. 

Let $Q_1(z) = \frac{P(z)}{1-z/\lambda_1}$ and let $Q(z) = \frac{Q_1(z)}{1-z/\lambda_2} = \frac{P(z)}{(1-z/\lambda_1)(1-z/\lambda_2)}$. Since $P(z)$ has roots at $\lambda_1$ and $\lambda_2$ and is holomorphic on $\mathbb{D}$ we know that $Q_1(z)$ and $Q(z)$ are both also holomorphic on $\mathbb{D}$. If we write out their power series expansions $Q_1(z) = \sum_{j=0}^\infty q_{1, j}z^j$ and $Q(z) = \sum_{j=0}^\infty q_jz^j$ then the $q_j$ they satisfy the 2-step recursion:
\begin{align*}
q_0 &= p_0 \\
q_1 &= p_1 + 2(\sqrt{M}+1)\cos(\theta)q_0 \\
q_{j+2} &= p_{j+2} + 2(\sqrt{M}+1)\cos(\theta)q_{j+1} - (\sqrt{M}+1)^2q_j
\end{align*}

where $\cos(\theta)$ is the argument of $\lambda_1$. Note that since $\lambda_1$ could be real it may be that $\cos(\theta) = 1$. Also, note that the escape threshold of $Q(z)$ is exactly 1.

Notice $p_0 = q_{0, 0} = q_0 = 1$ which is exactly equal to the escape threshold for $Q(z)$. This implies that $|q_j| = 1$ for every index $j$. Since $q_j$ is real this means $q_j \in \{-1, 1 \}$. If there is an index $j$ such that $q_j = 1$ and $q_{j+1} = -1$ then $|q_{j+2}|$ will be $> 1$ giving a contradiction, therefore since $q_0 = 1$ we have $q_j = 1$ for every $j$ and hence that $p_j = M$ for every $j \geqslant 2$. Thus $Q(z) = \sum_{j=0}^\infty z^j$ and hence $P(z) = (1-z/\lambda_1)(1-z/\lambda_2)\sum_{j=0}^\infty z^j$. On the other hand, we know that $p_0 = 1$ and $p_2 = p_3 = ... = M$. This is only possible if $\lambda_1 = \lambda_2 = \frac{1}{\sqrt{M}+1}$ and if $p_1 = -(2\sqrt{M} + 1)$ as claimed.



\bigskip
\noindent \textit{Proof of (2)}


Suppose $P_n(z) = \sum_{j=0}^\infty p_{n, j}z^j$ are power series with $p_{n, j} \in S$ and with roots $\lambda_{n, 1}$ and $\lambda_{n, 2}$ converging to $\frac{1}{\sqrt{M} + 1}$ and that $p_{n, 0} = 1$ and $p_{n, 1} = -(2\sqrt{M} + 1)$. Let $P(z) = 1-(2\sqrt{M}+1)z + \sum_{j=2}^\infty Mz^j$. Then we know that:
\begin{align*}
|P(\frac{1}{\sqrt{M} + 1}) - P_n(\frac{1}{\sqrt{M}+1}) | \rightarrow 0.
\end{align*}

Let $D_n(z) = P(z) - P_n(z) = \sum_{j=2}^\infty (M-p_{n, j})z^j$ Then we have $D_n(\frac{1}{\sqrt{M}+1}) \rightarrow 0$. Note that $D_n(z) > 0$ for every positive real number $z$.

If there is some index $j$ such that for infinitely many $n$, $p_{n, j} \neq M$ then $D_n(\frac{1}{\sqrt{M}+1}) \geqslant (M-p_{n, j})(\frac{1}{\sqrt{M}+1})^j$ which means infinitely many $D_n(\frac{1}{\sqrt{M}+1})$ have a uniform non-zero lower bound and hence that this can not converge to $0$.

Thus we must have that for every index $j$ there exists an $N$ such that $n > N$ implies for every $2 \leqslant k < n$, $p_{k, j} = M$.


\bigskip
\noindent \textit{Proof of (3)}

If $|p_{n, 0}| > 1$ then $P_n(z)$ can not have two different roots both deeper then $\frac{1}{\sqrt{M/p_0}+1}$ and hence they can not approach $\frac{1}{\sqrt{M}+1}$. Thus we can assume $|p_{n, 0}| = 1$ for $n$ large enough and hence, after possibly multiplying by $-1$, that $p_{n, 0} = 1$.

If $p_{n, 1} \neq -(2\sqrt{M}+1)$ for all $n$ then because $p_{n, 1} \in S$ and $S$ is finite we must have $|p_{n, 1} + (2\sqrt{M}+1)|$ is uniformly $> 0$.

if $p_{n, 1} < -(2\sqrt{M}+1)$ for infinitely many $n$ then for such $n$ and for any positive real number $z$, $P_n(\frac{1}{\sqrt{M}+1})$ is uniformly $< 1 -(2\sqrt{M}+1)\frac{1}{\sqrt{M}+1} + \sum_{j=2}^\infty M(\frac{1}{\sqrt{M}+1})^j = 0$ and hence $P_n(\frac{1}{\sqrt{M}+1})$ does not converge to $0$ which is a contradiction.

Now suppose $p_{n, 1} > -(2\sqrt{M}+1)$ for infinitely many $n$. Again this means that $2\sqrt{M}+1 + p_{n, 1}$ is uniformly $> 0$ for such $n$. Let $\lambda_n$ be a complex root of $P_n$ in the first quadrant such that $\lambda_n \rightarrow \frac{1}{\sqrt{M}+1}$. Then we have $\lambda_n = \frac{1}{r_n}e^{i\theta_n}$ where $\theta_n \rightarrow 0$ and $r_n \rightarrow \sqrt{M}+1$.

Consider the 2-step recursion:
\begin{align*}
p_{n, 0} &= q_{n, 0}\\
p_{n, 1} &= q_{n, 1} - 2r_n\cos(\theta_n)q_{n, 0} \\
p_{n, j+2} &= q_{n, j+2} - 2r_n\cos(\theta_n)q_{n, j+1} + r^2_nq_{n, j}
\end{align*}

We know that $p_{n, 0} = 1$ and hence $q_{n, 0} = 1$. Then $q_{n, 1} = p_{n, 1} + 2r_n\cos(\theta)$. Notice that as $n \rightarrow \infty$ we have $2r_n\cos(\theta) \rightarrow 2\sqrt{M} + 2$. On the other hand, $p_{n, 1}$ is uniformly $> -(2\sqrt{M} + 1)$ meaning $q_{n, 1}$ is uniformly $> 1$. But the escape threshold for the $q_{n, j}$ tends to $1$ as $n \rightarrow \infty$ giving a contradiction.

Since, for infinitely many $n$ having $p_{n, 1} < -(2\sqrt{M}+1)$ or $p_{n, 1} > -(2\sqrt{M}+1)$ are both impossible, we must have that for $n$ large enough, $p_{n, 1} = -(2\sqrt{M}+1)$ as wanted.

\end{proof}

\end{document}
